\documentclass[11pt]{article}
\usepackage[margin=1.1in]{geometry}
\usepackage{amsmath,amssymb,amsthm}
\usepackage{microtype}
\usepackage{booktabs}
\usepackage[hidelinks]{hyperref}
\usepackage{enumitem}

\newtheorem{definition}{Definition}
\newtheorem{proposition}{Proposition}

\title{Liquid Light and the Unsupported Nucleus:\\
\large A 2021 Private Cosmology as a Baseline for Generative Completion}
\author{Flyxion\\ \small Independent Researcher}
\date{September 2026}

\begin{document}
\maketitle

\begin{abstract}
A seventeen-page self-published ``theory of everything'' written during the 2021 lockdown, together with a condensed restatement posted publicly some months later, is examined as a specimen of private cosmology produced before conversational language models were widely available. The documents propose that light has a gaseous phase (visible light, identified with the wave description) and a liquid phase (dark, identified with the particle description), and that an inward flow of the latter causes gravity and constitutes dark matter. The analysis has three aims. First, it characterizes the method of the documents: analogy promoted to identity, drift of technical categories, and the absence of any second party. Second, it isolates the intuitions that have genuine counterparts in physics (push gravity, condensate dark matter, phases of the vacuum) and formalizes each only far enough to yield a number or a contradiction. The push-gravity reading, for instance, requires an absorbed power per kilogram of Earth exceeding the measured geothermal output by more than twenty orders of magnitude. Third, it uses the specimen as a baseline against which to describe what generative systems change. The central claim is that such systems do not create the appetite for total explanation; they remove the labour that once kept the unsupported inferential nucleus of a private cosmology visible, and they can manufacture apparent corroboration whose evidential weight is exactly zero. A criterion is proposed for distinguishing assisted theorizing from manufactured completion: an output counts as evidence only when it depends on something the premise did not generate.
\end{abstract}

\section{The Specimen}

Two documents are under consideration. The first, titled \emph{The Theory of Everything}, is a seventeen-page essay dated internally only ``Lockdown 2021'' and written in South Africa. It was deposited on Zenodo on 23 March 2021 under a CC-BY 4.0 licence \cite{orkadmon2021}, which fixes its date independently of the author's own dating. The second is a short public post from around September 2021, in which the author announces that he is redoing an earlier, messier post and addresses the proposal specifically to the problem of ``uncertainty'' in quantum physics. Both precede the public release of conversational language models in late 2022. The author disclaims formal training in physics and describes his method, in his own words, as ``common sense logic and a spooky jump into the unknown.''

The central inference of both documents can be stated in a sentence. Water vapour rises and appears wavelike; raindrops fall and appear particulate; therefore the wave description corresponds to a gas phase and the particle description to a liquid phase, and light, which admits both descriptions, must have both phases. Visible light is the gas. The liquid is dark, flows inward toward the centres of stars and planets, pushes baryonic matter along with it, and is what is experienced as gravity. It is also dark matter. Stars take it in, ``boil'' it, and emit it as visible light.

Two features make the pair useful for the purposes of this essay. The first is that the documents are dated and unassisted. Whatever their content, they record what a private cosmology looks like when every sentence has to be produced by the person who holds it. The second is that there are two of them, months apart. The later post is not a summary of the earlier essay alone. It adds material: the cosmic microwave background is reinterpreted as the motion of ``ground-state liquid light,'' the dark phase acquires a ``microwave black spectrum,'' and the target of the theory narrows from gravity and dark matter to the foundations of quantum mechanics. The growth between the two is therefore observable, and its rate and character are part of the evidence.

That growth is slow, local, and visibly the author's own. New claims arrive one at a time. None arrives with an equation, a citation to the literature it touches, or a response to an objection. The CMB reinterpretation, for example, is asserted without any engagement with the property that makes the CMB informative, namely that its spectrum is a blackbody at $2.725\,$K to a precision of tens of parts per million \cite{fixsen1996}. The later document is shorter, more confident, and no more complete.

A note on sources is required, since the essay's evidential weight rests partly on its fidelity to texts most readers will not consult. The first document is publicly archived with a persistent identifier, so every claim made about it here can be checked against the source; the file analysed is byte-identical to the archived deposit. Quotations from both documents are verbatim and kept short. The characterization of the documents is confined to claims each makes explicitly; nothing is attributed to the author that appears only in secondary summaries of his work, several of which circulate without sources. The September post is not archived in the same way, and readers who have not seen it have contact with it only through this essay.

\section{Method}

\subsection{Analogy promoted to identity}

The documents do not argue that gas and wave behaviour are analogous. They assert that ``a wave function rises up and a particle function falls down,'' and treat the observation of water as licensing the claim about light. Once the identity is in place, a large range of phenomena are read through it: birds of prey climbing and stooping, mosquitoes, photosynthesis with its light-dependent and light-independent reactions, lightning, the anatomy of the eye, dreams of falling, and the rhythm of day and night. Each is presented as further confirmation. None is independent of the original identification, because each is interpreted by means of it. The pattern is familiar from any sufficiently flexible interpretive scheme: once rising-and-expanding and falling-and-contracting are the only categories, almost anything that moves will fit one or the other.

\subsection{Category drift}

The dark phase of light is described, at different points, as a liquid, a solid, a superfluid, a type-II superconductor, a Bose--Einstein condensate, the source of entangled photons, a polar molecule, an exotic metal alloy, the gravitational boson, dark matter, and an electromagnetic field. It is said to be ``the heaviest densest substance of all while still remaining invisible and `massless'.'' These categories have definitions that are not interchangeable: a superconductor is characterized by zero resistance and flux expulsion in a charged system, a Bose--Einstein condensate by macroscopic occupation of a single mode, an alloy by a mixture of metallic elements. Moving among them without specifying the transitions or the observables that would distinguish them has a specific effect: no claim about the substance can be refuted, because any property demanded by an objection can be supplied by switching category.

\subsection{Correct premises, unlicensed inferences}

It would be inaccurate to describe the documents as uniformly mistaken. Several premises are correct and some are stated well. Hydrostatic equilibrium in stars is described accurately as a balance between inward pressure and outward radiation. Pascal's principle is cited correctly. The ``hydrogen wall'' at the edge of the heliosphere is a real structure, detected in Lyman-$\alpha$ absorption \cite{linsky1996} and later in New Horizons ultraviolet data. The failures lie in what is inferred from these facts. The hydrogen wall becomes the outer boundary of a gravity-generating bubble; hydrostatic equilibrium becomes evidence for an inward pushing fluid rather than a consequence of self-gravity. There are also direct factual errors that do not depend on any theoretical commitment: rods and cones are not located respectively in the pupil and the iris, and a cloud is not a single molecule.

\subsection{The solitary condition}

One sentence in the 2021 essay deserves more attention than its content. The author writes that it is surprising he cannot find the question anywhere but in his own head. This is the condition under which the documents were written: no interlocutor, no adjacent literature located, no second voice to elaborate or to resist. The author experiences this solitude as evidence of originality. It is better understood as the absence of friction and of corroboration alike. The documents contain no one to agree with them and no one to object, and their roughness is in large part the trace of that absence.

\section{The Seeds}

The documents contain intuitions that are not absurd, and which have real counterparts in the history or current practice of physics. A fair analysis should name them. The difference between the documents and their counterparts is not the idea; it is everything that turns an idea into a claim that could fail.

The seeds are not equally close to their counterparts, and treating them under one heading risks overstating the nearer ones' company. Table~\ref{tab:seeds} ranks them by distance. Only the first shares a mechanism with a developed theory. The second shares a substance and an image but none of the parameters. The third and fourth share little more than a phrase: the relevant programmes consist almost entirely of mathematical apparatus the documents do not approach, and in the case of vacuum phases the specific mapping the documents propose contradicts the physics it gestures at.

\begin{table}[h]
\centering\small
\begin{tabular}{@{}p{2.6cm}p{2.9cm}p{3.6cm}p{3.6cm}@{}}
\toprule
Seed & Counterpart & Shared & Missing \\
\midrule
Push gravity & Le Sage & The mechanism: shielding of a pervasive flux gives an inverse-square force & Cross-sections, flux intensity, and the energy budget, which is where it fails \\
Condensate dark sector & Berezhiani--Khoury & The substance: a superfluid in halos & Particle mass, coupling, transition temperature, rotation-curve predictions \\
Vacuum phases & Volovik; analogue gravity & The image of space as a medium & The entire apparatus (Lagrangian, order parameter, derived excitations); the wave/particle mapping contradicts complementarity \\
Thermodynamic gravity & Jacobson & A slogan: thermodynamics explains gravity & Which heat, temperature, and entropy; any derivation \\
\bottomrule
\end{tabular}
\caption{The seeds, ordered by distance from their developed counterparts.}
\label{tab:seeds}
\end{table}

\subsection{Gravity as a push}

The proposal that gravity is caused by a pervasive flux pushing matter together is the kinetic theory of gravitation associated with Nicolas Fatio and Georges-Louis Le Sage \cite{lesage1784}. In Le Sage's version, space is filled with an isotropic flux of small, fast ``ultramundane'' corpuscles. An isolated body is struck equally from all sides and feels no net force. Two bodies partially shield each other, so each receives less flux from the direction of the other and is pushed toward it. Because the shadow cast by one body on another subtends a solid angle proportional to $1/r^2$, the resulting force falls off as the inverse square of the distance. This is a genuine derivation, and it is the most interesting thing the 2021 documents half-contain. Appendix~\ref{app:lesage} works it out. The theory was abandoned in the nineteenth century for quantitative reasons, and those reasons apply with greater force to the liquid-light variant.

\subsection{The dark sector as a condensate}

The claim that dark matter is a superfluid condensate is not in itself eccentric. Berezhiani and Khoury have proposed that dark matter consists of light particles that form a superfluid in galactic halos, where the phonon excitations of the condensate mediate an additional force that reproduces the empirical regularities usually summarized as MOND, while the substance behaves as ordinary cold dark matter on cosmological scales \cite{berezhiani2015}. The proposal specifies a particle mass range, a transition temperature, an equation of state, and observational consequences for galaxy rotation curves and cluster dynamics. Photons can also be made to condense, but only under engineered conditions: in a dye-filled optical microcavity that confines photons, gives them an effective mass, and approximately conserves their number \cite{klaers2010}. Free thermal radiation does not condense, because photon number is not conserved. The 2021 documents gesture at the first program and implicitly contradict the conditions of the second.

\subsection{The vacuum as a medium with phases}

The mapping of wave and particle descriptions onto gas and liquid phases is mistaken, since wave and particle are complementary aspects of a single quantum state rather than two states of aggregation. But the broader thought, that what appears as empty space is a medium with phases, is mainstream. Spontaneous symmetry breaking treats the vacuum as having ordered and disordered phases. Volovik has developed a detailed programme in which the low-energy physics of the vacuum is modelled on superfluid helium \cite{volovik2003}, and the analogue-gravity literature studies how effective metrics arise for excitations in condensed-matter media \cite{barcelo2011}. These programmes share with the 2021 documents an image and very little else. They share with each other a discipline: the medium is specified by a Lagrangian or an equation of state, and the emergent physics is derived rather than asserted.

\subsection{Emergent gravity}

There is also a serious tradition that treats gravity as emergent from something more primitive. Jacobson derived the Einstein field equations from the Clausius relation $\delta Q = T\,dS$ applied to local causal horizons \cite{jacobson1995}. The 2021 essay invokes the first law of thermodynamics as the true explanation of stars and gravity, and in that respect it points at a real idea. What Jacobson's derivation has and the essay lacks is a statement of which heat, which temperature, and which entropy, and a derivation that returns a known equation as its output.

\section{Formalization as a Test, Not a Rescue}

At this point the analysis faces a methodological hazard that is the subject of the remainder of the essay. To formalize the seeds of the 2021 documents is to supply the connective tissue the documents lack: the cross-sections, the equations of state, the definitions. That is precisely the operation a capable language model performs when asked to develop such a theory. If it is done without discipline, the result is a rescued theory whose apparent coherence comes entirely from the rescuer.

The discipline adopted here has two rules. First, every assumption not supplied by the original author is marked as such. Second, each formalization is pushed only until it produces a number that can be checked against observation, or an internal contradiction, and then stopped. The aim is not to make the theory as good as possible. It is to find the point at which it makes contact with something it did not generate.

Applied to the push-gravity reading, this procedure terminates quickly. Appendix~\ref{app:lesage} shows that a flux able to reproduce Newtonian gravity while keeping gravitational force proportional to mass must be absorbed so weakly that, for Earth, the absorbed power exceeds the measured geothermal heat flow by a factor of order $10^{23}$. The same flux produces a drag on Earth's orbital motion more than a hundred times stronger than the Sun's gravitational attraction. These are the objections raised by Maxwell and Poincaré against Le Sage \cite{maxwell1875,poincare1908}, and they apply to any variant in which the pushing medium is absorbed. The liquid-light variant makes them worse. Its medium is said to move ``nowhere near'' the speed of light, which increases aberration; Laplace had already shown that an orbital force propagating at finite speed must propagate at millions of times the speed of light to avoid destabilizing the Moon's orbit \cite{laplace1805}.

Applied to the day--night claim, the procedure terminates immediately. The 2021 essay states that at night we experience a ``double'' gravitational effect. Superconducting gravimeters measure local gravity continuously with a resolution far better than one part in $10^{10}$ of $g$, and the diurnal variations they record are accounted for, in phase and amplitude, by solar and lunar tides of order $10^{-7}g$ \cite{hinderer2015}. A doubling would exceed the observed effect by roughly seven orders of magnitude. The prediction, once written as a number, is refuted by data that already exist.

Applied to the condensate reading, the procedure does not terminate in refutation, and that is informative too. A dark-matter condensate is an open possibility with published constraints. But none of the constraints on it bear on the 2021 documents, because the documents supply none of the parameters the constraints are about: no particle mass, no coupling, no transition temperature, no rotation curve. The documents are not wrong about condensate dark matter. They have not yet said anything about it.

\section{What Generative Systems Change}

\subsection{The production function}

The 2021 documents show that generative systems did not invent private cosmology. The appetite for total explanation, the joining of analogy, mystical correspondence, misunderstood terminology, and dissatisfaction with institutional science, was fully present before them. What changed is the cost of producing the external form of scholarship.

Before conversational models, a seventeen-page theory of everything required enough personal fixation to write seventeen pages, and the result usually carried visible traces of its origin: shifting definitions, unanswered objections, missing mathematics, repetition, typographical error. The 2021 essay ends with a scriptural line repeated eight times. These features were never evidence of error, and polish was never evidence of grounding. Both were proxies for labour, and labour was a proxy for commitment and some minimal competence. What generative systems break is the proxy chain. A person can now supply three metaphors and receive an abstract, numbered axioms, field equations, figures, a bibliography, proposed experiments, and replies to anticipated objections. The inferential nucleus is unchanged. Its surface distance from legitimate scholarship collapses.

\subsection{Completion without contact}

The operations a model performs on such material are specific. It supplies continuity between claims that had none. It resolves verbal contradictions (heaviest substance, yet massless) by choosing one reading and building on it, without resolving any physical contradiction. It invents intermediate mechanisms. It converts objections into further elaboration. Each operation removes a visible symptom of incompleteness. None adds a relation to the world that was not already present. The result may be described as \emph{manufactured completion}: the production of document-completeness around an unsupported nucleus.

The rate difference is the important historical fact. Between the lockdown essay and the September post, a period of months, the 2021 theory acquired a handful of new claims. A generative system could produce a monograph, a repository, a simulation, and a set of explanatory dialogues from the same nucleus within days. What formerly exhausted itself in a comment thread can now acquire the full institutional apparatus of a research programme.

This rate claim is a hypothesis, not a finding of this essay, and it is testable. If generative systems change the production function rather than the underlying conviction, then in venues that receive unsolicited theoretical submissions, the number of pages per author and the density of formal apparatus (equations, citations, numbered definitions) per submission should have risen after late 2022 considerably faster than the number of distinct authors. A rise in authors without a rise in apparatus density would count against the hypothesis.

\subsection{Corroboration without independence}

The more serious change concerns corroboration. The author of the 2021 essay could find no one else who held his view. A present-day author need not look far. Asked whether the theory has merit, a model will reconstruct the strongest version of it, locate the real adjacent literature (photon condensates, superfluid dark matter, emergent gravity), and may reasonably say that the intuition connects to serious work. That answer is then available to be cited as confirmation. A second model summarizes the first model's elaboration; a third supplies equations; a fourth evaluates the equations within assumptions inherited from the same premise. Apparent corroboration accumulates.

Appendix~\ref{app:witness} states the result formally. When every response is generated from the same premise, and the responses do not depend on any observation of the world beyond that premise, the responses are conditionally independent of the truth of the hypothesis given the premise, and their agreement has a likelihood ratio of exactly one. Under a weaker assumption, that the responses share a correlated error, the effective number of independent witnesses among $N$ is bounded by $1/\rho$ for correlation $\rho$, however large $N$ becomes. Agreement among such witnesses measures the shared prior. It does not measure the world.

This is the mechanism by which a private cosmology acquires the look of a literature. The danger is not that unsupported theories can be written; they always could. It is that completion, criticism, corroboration, and apparent peer response can all now be generated from the same premise without any new contact with the world.

\section{A Criterion}

An argument of this kind will be read reflexively. Any large body of theoretical work produced with the assistance of language models, and in particular any that proposes a unifying framework across physics, cognition, and other domains, invites the same diagnosis. The argument must therefore state what distinguishes legitimate assisted theorizing from manufactured completion, or it cannot exempt anything, including itself.

The distinction proposed here does not depend on who or what produced the text, nor on its surface quality. It depends on the dependency structure of its claims.

\begin{definition}[Contact]
An output $O$ makes \emph{contact} relative to a premise $P$ if $O$ depends on some variable $W$ that is not a function of $P$ and the generating process alone, and whose value could have made $O$ fail.
\end{definition}

Examples of $W$ include a measurement, a replicated experimental result, the output of a program run against inputs the premise did not choose, a proof checked by a mechanical verifier, or a published dataset. Examples of outputs without contact include a restatement of the premise in formal notation, an equation whose parameters are chosen to reproduce the premise, a second model's endorsement, and a list of adjacent literatures. These may be valuable as clarification. They are not evidence.

The criterion can be seen operating in miniature in the analysis of the 2021 documents. A first reading of the condensed September post produced a set of objections. A second reading, of the full lockdown essay, confirmed and extended them. Superficially this resembles the corroboration cascade described above. It differs in one respect: the second reading made contact with a document the first had not seen, and could have found that the condensed post misrepresented the full argument. It did not, and that absence of failure carries weight precisely because failure was possible.

The example requires a qualification. A document can serve as $W$, but not because documents are inherently contact; documents can be manufactured from a premise as easily as any other output, which is the whole difficulty. The lockdown essay qualifies because of its causal history: it predates the September post and was not generated from it or from the analysis of it. Contact is a property of the path by which $W$ came to be, not of the kind of object $W$ is. The same test applies to every example in the list above. A dataset produced by a pipeline tuned to the hypothesis, or a proof of a lemma chosen because it was provable, can fail the criterion while looking like a paradigm case of meeting it.

A second qualification concerns how much a passed test is worth. Contact is necessary for evidence but not sufficient for much of it. When an output survives a test, the likelihood ratio in its favour is $\Pr(\text{pass}\mid H)/\Pr(\text{pass}\mid\neg H)$, and a test that would usually be passed even if the hypothesis were false supports it only weakly. Contact therefore comes in degrees, which the philosophy of statistics has studied under the name of severity \cite{mayo1996}. The second reading of the 2021 material was a test of moderate severity: a condensed post misrepresenting its source is not rare, but neither is it the usual case. The calculations in the appendices are severe tests: had push gravity been viable, there was no reason for the heating rate to exceed the geothermal output by twenty-three orders of magnitude rather than falling below it. The same is true of the calculations in the appendices. The numbers could have come out otherwise; the geothermal flux and the gravimeter records are not functions of the premise.

The formalizations in Section~4 were written to satisfy this criterion rather than to rescue the theory, and the criterion explains why the procedure was stopped where it was. Once a formalization produces a number that the world contradicts, further elaboration can only restore coherence by adding assumptions whose sole function is to avoid the contradiction. That is manufactured completion by definition.

\section{Conclusion}

The 2021 documents are useful because they are rough. Their roughness exposes the inferential nucleus that a contemporary generative system would conceal beneath competent exposition: an analogy between rising vapour and falling rain, promoted to an identity and extended to light, gravity, dark matter, and the foundations of quantum mechanics. Within that nucleus there are genuine seeds, which physics has in several cases developed seriously: push gravity, condensate dark matter, phases of the vacuum, and emergent gravity. What separates the seeds from their developed counterparts is not imagination. It is the specification of quantities that could fail, and the willingness to let them.

Generative systems did not invent private cosmology. They automated its passage from intuition to literature. The appropriate response is not to distrust fluent text as such, since fluency was never the evidence, but to ask of every document, however produced, what in it could have turned out otherwise.

\appendix

\section{The Le Sage Force and Its Costs}
\label{app:lesage}

This appendix supplies assumptions not present in the 2021 documents, marked here as such, in order to extract numbers.

\paragraph{Assumption (supplied).} Space is filled with an isotropic flux of massless carriers, with momentum flux per unit area per unit solid angle $I$ (units of pressure per steradian). A body of mass $m$ absorbs the flux with effective cross-section $\sigma = k m$, where $k$ is a constant (units $\mathrm{m^2\,kg^{-1}}$). Absorption is complete and inelastic.

\paragraph{The force.} Consider bodies $A$ and $B$ separated by $r$, with $r$ large compared to their sizes. As seen from $A$, body $B$ subtends a solid angle $\Omega_B = \sigma_B / r^2$ and removes the flux from that direction. The net momentum delivered to $A$ per unit time, directed toward $B$, is therefore
\begin{equation}
F = I\,\sigma_A\,\Omega_B = \frac{I\,\sigma_A\,\sigma_B}{r^2} = \frac{I k^2\, m_A m_B}{r^2}.
\end{equation}
Matching to Newton's law requires
\begin{equation}
G = I k^2.
\end{equation}

\paragraph{Mass proportionality.} The linear relation $\sigma = km$ holds only if bodies are nearly transparent to the flux. For a body of column density $\Sigma$, the absorbed fraction is $1 - e^{-k\Sigma}$, which is linear in mass only when $k\Sigma \ll 1$. Otherwise the outer layers shadow the inner ones and gravitational force ceases to be proportional to mass, contrary to the equivalence of gravitational and inertial mass tested to high precision. For Earth, the column density through the centre is $\Sigma \approx 2R_\oplus \bar\rho \approx 7.0\times10^{10}\,\mathrm{kg\,m^{-2}}$. Requiring deviations from proportionality below one percent gives
\begin{equation}
k \lesssim \frac{0.01}{\Sigma} \approx 1.4\times10^{-13}\,\mathrm{m^2\,kg^{-1}}.
\end{equation}
This is a lenient bound; equivalence-principle tests would tighten it by many orders of magnitude, strengthening every conclusion below.

\paragraph{Heating.} Carriers arriving within solid angle $d\Omega$ deliver momentum at rate $I\sigma\,d\Omega$ in magnitude. For massless carriers, energy and momentum are related by $E = pc$, so the energy absorbed from that solid angle per unit time is $cI\sigma\,d\Omega$. Integrating over the full sphere, with the flux isotropic,
\begin{equation}
\dot W = \int_{4\pi} c\,I\sigma\,d\Omega = 4\pi c\,I\sigma.
\end{equation}
The directions cancel in the net force but not in the absorbed energy, which is why a body at rest feels nothing and is nonetheless heated. Using $I = G/k^2$ and $\sigma = km$,
\begin{equation}
\frac{\dot W}{m} = \frac{4\pi c\, G}{k} \gtrsim 1.8\times10^{12}\ \mathrm{W\,kg^{-1}}.
\end{equation}
Earth's measured surface heat flow is about $47\,$TW \cite{davies2010}, or about $7.9\times10^{-12}\,\mathrm{W\,kg^{-1}}$. The required absorption exceeds it by a factor of order $10^{23}$. This is Maxwell's and Poincaré's objection in quantitative form \cite{maxwell1875,poincare1908}. Elastic scattering avoids the heating but also removes the shadow, since scattered flux refills the shadowed directions; an inelastic flux that re-emits the energy in some other form must then explain where that form goes.

\paragraph{Drag.} Let the body move with velocity $\mathbf v$, $\beta = v/c \ll 1$, and let $\hat{\mathbf n}$ denote the direction from the body toward the source of arriving carriers, with $\cos\theta = \hat{\mathbf n}\cdot\hat{\mathbf v}$. In the body's rest frame, to first order in $\beta$, the specific intensity of a flux that is isotropic in the background frame becomes
\begin{equation}
I'(\theta) = I\,(1 + 4\beta\cos\theta),
\end{equation}
the factor of four collecting the Doppler shift of energy, the increased arrival rate, and aberration of solid angle, as for any massless isotropic bath. Carriers from direction $\hat{\mathbf n}$ push along $-\hat{\mathbf n}$, so the force is
\begin{equation}
\mathbf F = -\sigma\int_{4\pi} I'(\theta)\,\hat{\mathbf n}\,d\Omega
= -4\beta\,\sigma I\,\hat{\mathbf v}\int_{4\pi}\cos^2\theta\,d\Omega
= -\frac{16\pi}{3}\,\beta\,\sigma I\,\hat{\mathbf v},
\end{equation}
since the isotropic term integrates to zero and $\int\cos^2\theta\,d\Omega = 4\pi/3$. Writing $u = 4\pi I$ for the momentum-flux density recovers the familiar form $F = \tfrac43\sigma u\beta$ of radiative drag. Dividing by $m$ and using $I = G/k^2$, $\sigma = km$,
\begin{equation}
a_{\mathrm{drag}} = \frac{4}{3}\,\frac{4\pi G}{k}\,\frac{v}{c}.
\end{equation}
The result assumes complete absorption with no re-emission; any re-emission isotropic in the body frame carries no net momentum and leaves it unchanged at this order. For Earth's orbital speed, $v/c \approx 9.9\times10^{-5}$, and the bound on $k$ above, $a_{\mathrm{drag}} \gtrsim 0.78\,\mathrm{m\,s^{-2}}$, about $130$ times the Sun's gravitational acceleration at Earth ($5.9\times10^{-3}\,\mathrm{m\,s^{-2}}$). Earth's orbit would not survive a year.

\paragraph{The liquid-light variant.} The 2021 documents differ from Le Sage in two respects, both unfavourable. The flux is not isotropic but directed inward toward stars and planets, which requires a sink at every centre of mass whose capacity is nowhere specified; a planet, unlike a star, is not said to convert the inflow into anything. And the flux is slow, which enlarges the aberration that Laplace's analysis already constrained \cite{laplace1805}.

\section{Correlated Testimony}
\label{app:witness}

\paragraph{Conditional independence.} Let $H$ be a hypothesis about the world, $P$ a premise supplied by an author, and $O_1,\dots,O_N$ outputs of generative systems prompted with $P$. Suppose each output is produced as $O_i = f_i(P, \epsilon_i)$, where the noise terms $\epsilon_i$ are independent of $H$. Then $O = (O_1,\dots,O_N)$ is conditionally independent of $H$ given $P$, and
\begin{equation}
\frac{\Pr(O \mid H, P)}{\Pr(O \mid \neg H, P)} = 1.
\end{equation}
By Bayes' theorem, $\Pr(H \mid O, P) = \Pr(H \mid P)$. In the language of causal graphs, $P$ d-separates $H$ from every output \cite{pearl2009}. No number of outputs, and no degree of agreement among them, changes the credence warranted by the premise alone.

\begin{proposition}
If outputs are generated from a premise without dependence on any variable outside the premise and the generating process, their agreement carries no evidential weight regarding the hypothesis beyond that of the premise.
\end{proposition}

The Definition of contact in Section~6 is the negation of this hypothesis: contact is exactly what breaks the d-separation, by introducing a path from $H$ to $O$ through some $W$ that $P$ does not screen off.

\paragraph{Partial correlation.} In practice outputs may depend weakly on information outside the premise, for instance on training data that encodes real physics. A simpler model then applies. Suppose each of $N$ outputs carries an error of equal variance, with pairwise correlation $\rho$ arising from the shared premise and shared training. The effective number of independent witnesses is \cite{kish1965}
\begin{equation}
N_{\mathrm{eff}} = \frac{N}{1 + (N-1)\rho},
\end{equation}
which satisfies $N_{\mathrm{eff}} \to 1/\rho$ as $N\to\infty$. For $\rho = 0.9$, no number of outputs amounts to more than about $1.1$ independent witnesses. Adding models from the same family, or prompting one model repeatedly, raises $N$ without materially raising $N_{\mathrm{eff}}$.

\begin{thebibliography}{99}

\bibitem{barcelo2011}
C.~Barcel\'o, S.~Liberati, and M.~Visser, ``Analogue gravity,'' \emph{Living Reviews in Relativity} \textbf{14}, 3 (2011).

\bibitem{berezhiani2015}
L.~Berezhiani and J.~Khoury, ``Theory of dark matter superfluidity,'' \emph{Physical Review D} \textbf{92}, 103510 (2015).

\bibitem{davies2010}
J.~H.~Davies and D.~R.~Davies, ``Earth's surface heat flux,'' \emph{Solid Earth} \textbf{1}, 5--24 (2010).

\bibitem{fixsen1996}
D.~J.~Fixsen et al., ``The cosmic microwave background spectrum from the full COBE FIRAS data set,'' \emph{Astrophysical Journal} \textbf{473}, 576 (1996).

\bibitem{hinderer2015}
J.~Hinderer, D.~Crossley, and R.~J.~Warburton, ``Superconducting gravimetry,'' in \emph{Treatise on Geophysics}, 2nd ed., vol.~3 (Elsevier, 2015).

\bibitem{jacobson1995}
T.~Jacobson, ``Thermodynamics of spacetime: the Einstein equation of state,'' \emph{Physical Review Letters} \textbf{75}, 1260 (1995).

\bibitem{kish1965}
L.~Kish, \emph{Survey Sampling} (Wiley, 1965).

\bibitem{klaers2010}
J.~Klaers, J.~Schmitt, F.~Vewinger, and M.~Weitz, ``Bose--Einstein condensation of photons in an optical microcavity,'' \emph{Nature} \textbf{468}, 545--548 (2010).

\bibitem{laplace1805}
P.-S.~Laplace, \emph{Trait\'e de m\'ecanique c\'eleste}, vol.~4 (1805).

\bibitem{lesage1784}
G.-L.~Le Sage, ``Lucr\`ece Newtonien,'' \emph{Nouveaux M\'emoires de l'Acad\'emie Royale des Sciences et Belles-Lettres de Berlin} (1784).

\bibitem{linsky1996}
J.~L.~Linsky and B.~E.~Wood, ``The $\alpha$ Centauri line of sight: D/H ratio, physical properties of local interstellar gas, and measurement of heated hydrogen (the `hydrogen wall') near the heliopause,'' \emph{Astrophysical Journal} \textbf{463}, 254 (1996).

\bibitem{mayo1996}
D.~G.~Mayo, \emph{Error and the Growth of Experimental Knowledge} (University of Chicago Press, 1996).

\bibitem{maxwell1875}
J.~C.~Maxwell, ``Atom,'' \emph{Encyclop\ae dia Britannica}, 9th ed. (1875).

\bibitem{orkadmon2021}
J.~Orkadmon, ``Theory of Everything: Gravity and Dark Matter Explained,'' Zenodo (23 March 2021), \href{https://doi.org/10.5281/zenodo.4630378}{doi:10.5281/zenodo.4630378}.

\bibitem{pearl2009}
J.~Pearl, \emph{Causality: Models, Reasoning, and Inference}, 2nd ed. (Cambridge University Press, 2009).

\bibitem{poincare1908}
H.~Poincar\'e, \emph{Science et m\'ethode} (Flammarion, 1908).

\bibitem{volovik2003}
G.~E.~Volovik, \emph{The Universe in a Helium Droplet} (Oxford University Press, 2003).

\end{thebibliography}

\end{document}
